\documentclass[12pt,a4paper]{ctexart}

\usepackage{amsmath,amssymb}
\usepackage{geometry}
\usepackage{enumitem}
\usepackage{hyperref}
\usepackage{booktabs}
\usepackage{listings}

\geometry{left=2.5cm,right=2.5cm,top=2.5cm,bottom=2.5cm}
\setlist{nosep,leftmargin=2em}

\title{第7章\quad 机械振动}
\author{}
\date{}

\begin{document}

\maketitle

\section{知识点总结}

\subsection{机械振动基本概念}

\begin{itemize}
    \item \textbf{机械振动}：物体在一定位置附近作来回往复的运动。
    \item \textbf{广义振动}：任一物理量（位移、电流等）在某一数值附近反复变化。
    \item 振动形式：机械振动、电磁振动等。
\end{itemize}

\subsection{简谐振动（SHM）}

\begin{itemize}
    \item 物体离开平衡位置的位移 $x$ 随时间 $t$ 按余弦（或正弦）函数变化。
    \item \textbf{动力学特征}：所受合力的方向与位移方向相反，大小与位移成正比。
    \item \textbf{运动学特征}：位移满足二阶常系数线性齐次微分方程。
    \item \textbf{特点}：等幅振动、周期振动。
    \item \textbf{判断简谐振动的三要素}（同时满足）：
    \begin{enumerate}
        \item 受力特征 $F=-kx$；
        \item 运动微分方程形式 $\dfrac{\mathrm{d}^2x}{\mathrm{d}t^2}+\omega^2 x=0$；
        \item 运动方程形式 $x=A\cos(\omega t+\varphi)$。
    \end{enumerate}
\end{itemize}

\subsection{描述简谐振动的特征量}

\begin{itemize}
    \item \textbf{振幅 $A$}：最大位移的绝对值（参考圆的半径），由初始条件决定。
    \item \textbf{周期 $T$}：完成一次全振动所需的时间。
    \item \textbf{频率 $\nu$}：单位时间内完成全振动的次数。
    \item \textbf{角频率（圆频率）$\omega$}：旋转矢量的角速度。
    \item \textbf{相位 $\omega t+\varphi$}：决定振动物体在任一时刻的位置和运动状态。
    \item \textbf{初相 $\varphi$}：$t=0$ 时的相位，由初始条件决定。
    \item 三者关系：$\omega=2\pi\nu=\dfrac{2\pi}{T}$。
\end{itemize}

\subsection{由初始条件确定振幅与初相}

初始条件：$x_0=A\cos\varphi$，$v_0=-\omega A\sin\varphi$。

\begin{itemize}
    \item 振幅：$A=\sqrt{x_0^2+\dfrac{v_0^2}{\omega^2}}$。
    \item 初相：$\tan\varphi=-\dfrac{v_0}{\omega x_0}$（$\varphi$ 在 $-\pi\sim\pi$ 之间，需由初条件判断取舍）。
\end{itemize}

\subsection{速度与加速度}

\begin{itemize}
    \item 速度：$v=\dfrac{\mathrm{d}x}{\mathrm{d}t}=-\omega A\sin(\omega t+\varphi)$，$v_{\max}=\omega A$。
    \item 加速度：$a=\dfrac{\mathrm{d}^2x}{\mathrm{d}t^2}=-\omega^2 A\cos(\omega t+\varphi)$，$a_{\max}=\omega^2 A$。
    \item 相位关系：$v$ 与 $x$ 相差 $\pi/2$，$a$ 与 $v$ 相差 $\pi/2$，$a$ 与 $x$ 相差 $\pi$。
\end{itemize}

\subsection{旋转矢量表示法}

\begin{itemize}
    \item 自原点作大小为 $A$ 的矢量，以角速度 $\omega$ 逆时针旋转。
    \item 初始时与 $x$ 轴夹角为 $\varphi$，$t$ 时刻与 $x$ 轴夹角为 $\omega t+\varphi$。
    \item 矢端在 $x$ 轴的投影即简谐振动的位移 $x=A\cos(\omega t+\varphi)$。
\end{itemize}

\subsection{弹簧振子（固有频率）}

\begin{itemize}
    \item 角频率：$\omega=\sqrt{\dfrac{k}{m}}$。
    \item 周期：$T=2\pi\sqrt{\dfrac{m}{k}}$。
    \item 频率：$\nu=\dfrac{1}{2\pi}\sqrt{\dfrac{k}{m}}$。
    \item 周期和频率由振动系统决定；振幅、初相由初条件决定。
    \item \textbf{重要结论}：重力（或任一恒力）只影响谐振子系统的平衡位置，不影响固有频率。
\end{itemize}

\subsection{单摆（小角度 $\theta<5^\circ$ 近似下为谐振动）}

\begin{itemize}
    \item 动力学方程：$\dfrac{\mathrm{d}^2\theta}{\mathrm{d}t^2}+\dfrac{g}{l}\theta=0$。
    \item 角频率：$\omega=\sqrt{\dfrac{g}{l}}$。
    \item 周期：$T=2\pi\sqrt{\dfrac{l}{g}}$。
    \item 单摆周期与摆锤质量无关；可由 $T$、$l$ 测当地重力加速度 $g$。
\end{itemize}

\subsection{复摆（物理摆，小角度近似下为谐振动）}

\begin{itemize}
    \item 角频率：$\omega=\sqrt{\dfrac{mgl}{J}}$。
    \item 周期：$T=2\pi\sqrt{\dfrac{J}{mgl}}$。
    \item $J$ 为刚体对转轴的转动惯量，可由 $T$ 测定 $J$。
\end{itemize}

\subsection{简谐振动的能量}

\begin{itemize}
    \item 动能：$E_k=\dfrac{1}{2}mv^2=\dfrac{1}{2}kA^2\sin^2(\omega t+\varphi)$。
    \item 势能：$E_p=\dfrac{1}{2}kx^2=\dfrac{1}{2}kA^2\cos^2(\omega t+\varphi)$。
    \item 总能：$E=E_k+E_p=\dfrac{1}{2}kA^2$（与振幅平方成正比且守恒）。
    \item 周期平均值：$\overline{E_k}=\overline{E_p}=\dfrac{1}{2}E$。
\end{itemize}

\subsection{简谐振动的合成}

\subsubsection*{(1) 同方向、同频率合成}

\begin{itemize}
    \item 合振动仍为同频率的简谐振动 $x=A\cos(\omega t+\varphi)$。
    \item 振幅：$A=\sqrt{A_1^2+A_2^2+2A_1A_2\cos(\varphi_2-\varphi_1)}$。
    \item 初相：$\tan\varphi=\dfrac{A_1\sin\varphi_1+A_2\sin\varphi_2}{A_1\cos\varphi_1+A_2\cos\varphi_2}$。
    \item 同相（$\Delta\varphi=2k\pi$）时 $A=A_1+A_2$ 最大；反相（$\Delta\varphi=(2k+1)\pi$）时 $A=|A_1-A_2|$ 最小。
    \item $\Delta\varphi>0$ 时 $x_2$ 比 $x_1$ 超前。
\end{itemize}

\subsubsection*{(2) 同方向、不同频率合成——拍}

\begin{itemize}
    \item 合振动 $x=2A\cos\!\left(2\pi\dfrac{\nu_2-\nu_1}{2}t\right)\cos\!\left(2\pi\dfrac{\nu_2+\nu_1}{2}t\right)$。
    \item 拍频：$\nu_{\text{拍}}=|\nu_2-\nu_1|$。
    \item 拍周期：$T_{\text{拍}}=\dfrac{1}{|\nu_2-\nu_1|}$。
\end{itemize}

\subsubsection*{(3) 相互垂直的同频率合成}

\begin{itemize}
    \item 轨迹方程：$\dfrac{x^2}{A^2}+\dfrac{y^2}{B^2}-\dfrac{2xy}{AB}\cos(\beta-\alpha)=\sin^2(\beta-\alpha)$（椭圆类）。
    \item 特殊情形：
    \begin{itemize}
        \item 同相（$\beta-\alpha=0$）：$y=\dfrac{B}{A}x$（直线）；
        \item 反相（$\beta-\alpha=\pi$）：$y=-\dfrac{B}{A}x$（直线）；
        \item $\beta-\alpha=\dfrac{\pi}{2}$：$\dfrac{x^2}{A^2}+\dfrac{y^2}{B^2}=1$（顺时针椭圆，$A=B$ 时为圆）；
        \item $\beta-\alpha=\dfrac{3\pi}{2}$：$\dfrac{x^2}{A^2}+\dfrac{y^2}{B^2}=1$（逆时针椭圆，$A=B$ 时为圆）。
    \end{itemize}
\end{itemize}

\subsubsection*{(4) 相互垂直的不同频率合成}

\begin{itemize}
    \item 频率比为简单整数比时，轨迹为稳定的封闭曲线——李萨如图形，可用于测频率和相位差。
\end{itemize}

\subsection{阻尼振动}

\begin{itemize}
    \item 阻尼振动：振幅（或能量）随时间减小的振动。
    \item 阻尼因子：$\beta=\dfrac{\gamma}{2m}$；固有角频率：$\omega_0=\sqrt{\dfrac{k}{m}}$。
    \item 运动微分方程：$\dfrac{\mathrm{d}^2x}{\mathrm{d}t^2}+2\beta\dfrac{\mathrm{d}x}{\mathrm{d}t}+\omega_0^2 x=0$。
    \item \textbf{弱阻尼}（$\beta<\omega_0$）：$x=A_0 e^{-\beta t}\cos(\omega t+\varphi)$，$\omega=\sqrt{\omega_0^2-\beta^2}$，$T=\dfrac{2\pi}{\sqrt{\omega_0^2-\beta^2}}$。
    \item 振幅随时间指数衰减；周期比固有周期长。
    \item \textbf{临界阻尼}（$\beta=\omega_0$）与 \textbf{过阻尼}（$\beta>\omega_0$）：为非周期运动，缓慢回到平衡位置。
\end{itemize}

\subsection{受迫振动与共振}

\begin{itemize}
    \item 受迫振动：系统在周期性外力持续作用下所发生的稳定振动，频率等于驱动力的频率。
    \item 运动微分方程：$\dfrac{\mathrm{d}^2x}{\mathrm{d}t^2}+2\beta\dfrac{\mathrm{d}x}{\mathrm{d}t}+\omega_0^2 x=h\cos\omega t$。
    \item 稳定后：$x=A\cos(\omega t+\varphi)$，$A=\dfrac{h}{\sqrt{(\omega_0^2-\omega^2)^2+4\beta^2\omega^2}}$。
    \item \textbf{共振}：振幅达极大值的现象。
    \begin{itemize}
        \item 共振角频率：$\omega_{\text{共}}=\sqrt{\omega_0^2-2\beta^2}$。
        \item 共振振幅：$A_{\text{共}}=\dfrac{h}{2\beta\sqrt{\omega_0^2-\beta^2}}$。
        \item $\beta$ 越小，$\omega_{\text{共}}$ 越接近 $\omega_0$，$A_{\text{共}}$ 越大；$\beta=0$ 时 $A_{\text{共}}\to\infty$（尖锐共振）。
    \end{itemize}
    \item 防止共振：破坏外力周期性、改变系统固有频率、改变外力频率、增大阻尼。
\end{itemize}

\section{公式汇总}

\begin{enumerate}[leftmargin=3em,label=\textbf{公式\arabic*}：,ref=公式\arabic*]

    \item \textbf{简谐振动方程}：$x=A\cos(\omega t+\varphi)$

    说明：$A$ 振幅，$\omega$ 角频率，$\varphi$ 初相；由振动系统及初始条件决定。

    \item \textbf{动力学方程}：$F=-kx$

    说明：$k$ 劲度系数，$x$ 偏离平衡位置的位移；弹性恢复力与位移成正比反向。

    \item \textbf{运动微分方程}：$\dfrac{\mathrm{d}^2x}{\mathrm{d}t^2}+\omega^2 x=0$

    说明：$\omega^2=\dfrac{k}{m}$，是简谐振动的判别方程。

    \item \textbf{周期、频率、角频率关系}：$T=\dfrac{2\pi}{\omega}$，$\nu=\dfrac{1}{T}=\dfrac{\omega}{2\pi}$

    说明：描述振动快慢的三种等价参数。

    \item \textbf{弹簧振子}：$\omega=\sqrt{\dfrac{k}{m}}$，$T=2\pi\sqrt{\dfrac{m}{k}}$

    说明：$k$ 弹簧劲度系数，$m$ 振子质量；系统固有属性，与振幅无关。

    \item \textbf{振幅与初相}（初始条件）：$A=\sqrt{x_0^2+\dfrac{v_0^2}{\omega^2}}$，$\tan\varphi=-\dfrac{v_0}{\omega x_0}$

    说明：$x_0$、$v_0$ 为 $t=0$ 时的位移和速度。

    \item \textbf{速度、加速度}：$v=-\omega A\sin(\omega t+\varphi)$，$a=-\omega^2 A\cos(\omega t+\varphi)$

    说明：$v$ 比 $x$ 超前 $\pi/2$，$a$ 比 $v$ 超前 $\pi/2$（即 $a$ 与 $x$ 反相）。

    \item \textbf{单摆周期}：$T=2\pi\sqrt{\dfrac{l}{g}}$

    说明：$l$ 摆长，$g$ 重力加速度；仅在小角度（$\theta<5^\circ$）下成立。

    \item \textbf{复摆周期}：$T=2\pi\sqrt{\dfrac{J}{mgl}}$

    说明：$J$ 刚体对转轴的转动惯量，$m$ 刚体质量，$l$ 质心到转轴距离；小角度近似下成立。

    \item \textbf{简谐振动能量}：$E_k=\dfrac{1}{2}kA^2\sin^2(\omega t+\varphi)$，$E_p=\dfrac{1}{2}kA^2\cos^2(\omega t+\varphi)$，$E=\dfrac{1}{2}kA^2$

    说明：总能恒定且与振幅平方成正比；周期平均 $\overline{E_k}=\overline{E_p}=\dfrac{1}{4}kA^2$。

    \item \textbf{同方向同频率合成}：$A=\sqrt{A_1^2+A_2^2+2A_1A_2\cos(\varphi_2-\varphi_1)}$，$\tan\varphi=\dfrac{A_1\sin\varphi_1+A_2\sin\varphi_2}{A_1\cos\varphi_1+A_2\cos\varphi_2}$

    说明：合振动仍为同频率简谐振动；$\Delta\varphi=2k\pi$ 同相（振幅最大），$\Delta\varphi=(2k+1)\pi$ 反相（振幅最小）。

    \item \textbf{拍频}：$\nu_{\text{拍}}=|\nu_2-\nu_1|$

    说明：两同方向、不同频率且频率差很小的谐振动合成时，振幅周期性变化。

    \item \textbf{相互垂直同频率合成轨迹}：$\dfrac{x^2}{A^2}+\dfrac{y^2}{B^2}-\dfrac{2xy}{AB}\cos(\beta-\alpha)=\sin^2(\beta-\alpha)$

    说明：$\alpha$、$\beta$ 分别为 $x$、$y$ 方向振动的初相；轨迹为椭圆类二次曲线，相位差不同形状不同。

    \item \textbf{阻尼振动方程}（弱阻尼 $\beta<\omega_0$）：$x=A_0 e^{-\beta t}\cos(\omega t+\varphi)$，$\omega=\sqrt{\omega_0^2-\beta^2}$

    说明：$\beta=\dfrac{\gamma}{2m}$ 为阻尼因子，$\gamma$ 为阻力系数；振幅指数衰减。

    \item \textbf{受迫振动振幅}：$A=\dfrac{h}{\sqrt{(\omega_0^2-\omega^2)^2+4\beta^2\omega^2}}$

    说明：$h=H/m$ 为驱动力幅除以质量，$\omega$ 为驱动力角频率；稳态振幅。

    \item \textbf{共振条件}：$\omega_{\text{共}}=\sqrt{\omega_0^2-2\beta^2}$，$A_{\text{共}}=\dfrac{h}{2\beta\sqrt{\omega_0^2-\beta^2}}$

    说明：驱动力的角频率等于共振频率时，振幅达极大值。

\end{enumerate}

\end{document}
